\documentclass[a4paper,11pt]{article}
\usepackage{jheppub} % for details on the use of the package, please see the JINST-author-manual
\usepackage{lineno}
%\linenumbers
\usepackage{amsmath}
\DeclareMathOperator{\Tr}{Tr}

\arxivnumber{1234.56789} % if you have one

\title{\boldmath A title with some math: $x=1$}

% Collaborations

%% [A] If main author
%% \collaboration{\includegraphics[height=17mm]{collabroation-logo}\\[6pt]
%%  XXX collaboration}

%% or
%% [B] If "on behalf of"
%% \collaboration[c]{on behalf of XXX collaboration}


% Authors
% The "\note" macro will give a warning: "Ignoring empty anchor...", you can safely ignore it.

%% [A] simple case: 2 authors, same institution
%% \author[1]{A. Uthor\note{Corresponding author.}}
%% \author{and A. Nother Author}
%% \affiliation{Institution,\\Address, Country}

%% or, e.g.
%% [B] more complex case: 4 authors, 3 institutions, 2 footnotes
%% \author[a,b]{F. Irst,\note{Now at another university}}
%% \author[c]{S. Econd,}
%% \author[a,2]{T. Hird\note{Also at Some University.}}
%% \author[c,2]{and Fourth}
%% \affiliation[a]{Institution_1,\\Address, Country}
%% \affiliation[b]{Institution_2,\\Address, Country}
%% \affiliation[c]{Institution_3,\\Address, Country}

\author{A. Uthor}
\affiliation{One University,\\
some-street, Country}
\affiliation{Another University,\\
different-address, Country}

% E-mail addresses: only for the corresponding author
\emailAdd{first@one.univ}

\abstract{Abstract...}



\begin{document}
\maketitle
\flushbottom

\section{Introduction}
\label{sec:Introduction}

\section{Toroidal Near-Extremal Black Brane}
\label{sec:Toroidal Near-Extremal Black Brane}
In this section, we discuss the Toroidal Near-Extremal Black Brane solution in $AdS_4$ and the physically relevant scales that arise in this setup.\\

We start with the action,
\begin{equation}
    S=\frac{1}{16\pi G} \int_{ M} d^4x \sqrt{-g}(R-2\Lambda)-\frac{1}{4G}\int_{ M} d^4x \sqrt{-g}F^2 +\frac{1}{8\pi G}\int_{\partial M} d^3x \sqrt{-\gamma}K
\end{equation}

Which admits the following 4D metric as solution\\
\begin{equation}
    ds^2= -a(r)^2 dt^2 + \frac{dr^2}{a(r)^2} + b(r)^2(dx^2 +dy^2)
\end{equation}

\begin{align}
    a(r)^2=1-\frac{2GM}{r} + \frac{4\pi Q^2}{r^2} +\frac{r^2}{L^2}\\
    b(r)^2=r^2\\
    F_{xy}=\frac{Q_M}{r^2}
\end{align}
This is the charged Black brane solution with $M$ and $Q_M$ being the mass and magnetic charge of the Black brane respectively. $AdS_4$ radius L is related to $\Lambda$ by,
\begin{equation}
    L=\sqrt\frac{3}{|\Lambda|}
\end{equation}
The black brane solution becomes extremal at,\\
\begin{align}
    Q_{ext}=\frac{1}{4\pi}\big(r_h^2 +\frac{3r_h^4}{L^2}\big)\\
    M_{ext}=\frac{r_h}{G}\big(1 +\frac{2r_h^2}{L^2}\big)\\
    a^2(r)= \frac{(r-r_h)^2}{r^2L^2}(L^2 +3r_h^2+ 2rr_h+r^2)\\
\end{align}
Where $r_h$ is the horizon radius.\\    
We are interested in near-extremal black branes, therefore $M>Q$. The two zeroes of $a(r)$ are not co-incident but are at $r_+$ and $r_-$. Since we are close to extremality, 
\begin{align}
    r_\pm \approx r_h \pm \delta r_h \\
    \delta r_h \ll r_h \label{Eq: Near ext}
\end{align}

We work with large black branes, 
\begin{equation}\label{Eq: Large}
    r_+>>L 
\end{equation}

Then the near horizon region $r-r_+<<r_+$, has the geometry
\begin{equation}
    ds^2= -\frac{(r-r_+)(r-r_-)}{L_2^2}dt^2 + \frac{L_2^2}{(r-r_+)(r-r_-)}dr^2 +r_h^2(dx^2 +dy^2) + \mathcal{O}(\frac{r-r_+}{r_+})\ \text{corrections}
\end{equation}

So, the near horizon region factorizes into $AdS_2 \times S^2$, where
\begin{align}
    L_2 \approx \frac{L}{\sqrt{6}}\\
    R_{S^2} \approx r_h
\end{align}

The asymptotic $AdS_2$ boundary, where the $AdS_2$ throat connects to the $AdS_4$, is given by 
\begin{equation}
    r-r_h>>L_2
\end{equation}

Note that for a large black branes \ref{Eq: Large}, asymptotic $AdS_2$ boundary can lie in near horizon region,
\begin{equation}
    L_2 \ll r-r_h \ll r_h
\end{equation}

The chemical potential and temperature for the black brane is given by
\begin{align}
    \mu \sim \frac{r_h}{L^2} \\
    T\sim \frac{\delta r_h}{L^2}
\end{align}

Since we are close to extremality \ref{Eq: Near ext},
\begin{equation}\label{Eq:T\mu}
    \delta r_h \ll r_h \implies T \ll \mu
\end{equation}

We restate the chain of logic that sets the various scales. First, near-extremality sets $ \delta r_h \ll r_h$, then we work with large black branes which gives us $L\ll r_h$. Requiring that asymptotic $AdS_2$ region lies in near horizon region gives $L_2\ll r-r_h \ll r_h$.
Above conditions set $\mu \ll T$

\section{Kaluza- Klein reduction of $AdS_4$}
(Write why showing shear mode behaves like scalar helps us predict the leading 3rd order non linear correction in fluid gravity).\\
In previous sections, we analysed the four point function of a minimally coupled probe scalar in the near-extremal black brane background. While it is known that a minimally coupled scalar behaves same as the shear mode of metric peturbation in AdS backgrounds(at least at the 2 point function level), we show that it continues to hold for leading corrections to four point function in the Near extremal AdS background.\\

We start with the 4D action
\begin{equation}
S_4=\frac{1}{16\pi G_4}\int_{\mathcal M} d^4x\,\sqrt{-G}\,\big(R_{(4)}-2\Lambda\big)
+\frac{1}{8\pi G_4}\int_{\partial\mathcal M} d^3x\,\sqrt{-\Gamma}\,K
-\frac{1}{4G_4}\int d^4x\,\sqrt{-G}\,F^2 .
\label{eq:KK-S4}
\end{equation}
Using the KK ansatz,
\begin{equation}
G_{AB}\,dX^A dX^B=g_{\mu\nu}\,dx^\mu dx^\nu+\Phi_{ab}\,dy^a dy^b
\label{eq:KK-ansatz}
\end{equation}
(we work in $l=0$, hence KK gauge fields are absent).

The reduced action becomes:
\begin{equation}
\begin{aligned}
S={}&\frac{V_2}{16\pi G_4}\int_{\mathcal M} d^2x\,\sqrt{-g}\,\sqrt{+\det\Phi}\,
\Big[R-2\Lambda
+\underbrace{\frac14\,g^{\mu\nu}\,\text{Tr}\!\big(\Phi^{-1}\partial_\mu\Phi\big)\,\text{Tr}\!\big(\Phi^{-1}\partial_\nu\Phi\big)}_{\text{I}}\\
&\hspace{4.2cm}-\underbrace{\frac14\,g^{\mu\nu}\,\text{Tr}\!\big(\Phi^{-1}\partial_\mu\Phi\,\Phi^{-1}\partial_\nu\Phi\big)}_{\text{II}}\Big]\\
&+\frac{V_2}{8\pi G_4}\int_{\partial\mathcal M} dx\,\sqrt{-\gamma}\,\big(\det\Phi\big)^{1/2}K\;+\;S_{\text{Maxwell}}
\end{aligned}
\label{eq:KK-reduced}
\end{equation}

\textbf{Term I:}
using $\text{Tr}(\Phi^{-1}\partial_\mu \Phi)= \partial_\mu \text{Tr}(\text{ln} \Phi)= \partial_\mu \text{ln} \text{ det}(\Phi)$, let 
\begin{equation}
  \Phi= vM  \ \  \text{where} \ \ \det(M)=1 \implies  \det \Phi= v^2
\end{equation}
\begin{equation}
\text{I}=\frac14\,g^{\mu\nu}\,\partial_\mu\ln(v^2)\,\partial_\nu\ln(v^2) = \frac{1}{ v^2} \partial_\mu v \partial_\nu v
\label{eq:KK-termI}
\end{equation}

\textbf{Tern II:}\\
Substituting $\Phi=vM$, we get
\begin{equation}\label{Eq:term2}
\text{II}=-\frac14\,g^{\mu\nu}\,\text{Tr}\!\big(\Phi^{-1}\partial_\mu\Phi\,\Phi^{-1}\partial_\nu\Phi\big)=-\frac{g^{\mu \nu}}{2 v^2} \partial_\mu v \partial_\nu v - \frac{g^{\mu \nu}}{4} \text{Tr}(M^{-1}\partial_\mu M M^{-1}\partial_\nu M)
\end{equation}
Note that $\text{Tr}(M^{-1}\partial_\mu M)= \partial_\mu \text{ln} \det M=0$.\\
%=====================================================================
\\
\textbf{Maxwell term}

Now, we examine the Maxwell term,
\begin{equation}
\begin{aligned}
S_{Maxwell}=&-\frac{1}{4G_4}\int d^4x\,\sqrt{-G}\;F_{AB}F_{CD}\,G^{AC}G^{BD}\\
&\quad=-\frac{1}{4G_4}\int d^4x\,\sqrt{-g}\,\sqrt{\det\Phi}\,\Big(F_{\mu\nu}F_{\rho\sigma}\,g^{\mu\rho}g^{\nu\sigma}
+F_{ab}F_{cd}\,\Phi^{ac}\Phi^{bd}
+2\,F_{a\mu}F_{b\nu}\,\Phi^{ab}g^{\mu\nu}\Big).
\end{aligned}
\label{eq:KK-F2-reduce}
\end{equation}
Now, since all fields depend only on $(r,t)$,
\begin{equation}
\quad F_{xx}=F_{yy}=0,\quad F_{a\mu}=-\partial_\mu A_a .
\label{eq:KK-Fcomp}
\end{equation}
Therefore,
\begin{equation}
S_{\text{Maxwell}}=-\frac{V_2}{4G}\int d^2x\,\sqrt{-g}\,\sqrt{\det\Phi}\,\Big(F_{\mu\nu}F^{\mu\nu}
+2\,(\partial_\mu A_a)(\partial_\nu A_b)\,\Phi^{ab}g^{\mu\nu} + F_{xy}F_{xy} (\Phi^{xx} \Phi^{yy} - (\Phi^{xy}) ^2 )\Big)
\label{eq:KK-Smax-jordan}
\end{equation}
%Here, we can set  substitute $F_{xy}=Q_M$. \\

Combining all the terms, we get the KK reduced action
\begin{equation}
\begin{aligned}
S_2=&\frac{V_2}{16\pi G_4}\int_{\mathcal M} d^2x\,\sqrt{-g}\,v\,
\Big[R-2\Lambda + \frac{g^{\mu \nu}}{2 v^2} \partial_\mu v \partial_\nu v
- \frac{g^{\mu \nu}}{4}\text{Tr}(M^{-1}\partial_\mu M M^{-1}\partial_\nu M)\Big]\\
 & +\frac{V_2}{8\pi G_4}\int_{\partial\mathcal M} dx\,\sqrt{-\gamma}\,vK\;+ S_{Maxwell};
\end{aligned}
\label{eq:KK-reduced}
\end{equation}
\\
%=====================================================================
\textbf{Weyl rescaling to the Einstein frame}

Now, we Weyl rescale to go to the Einstein frame. A general Weyl rescaling in 2D gives
\begin{equation}
g\to\alpha(x)\,g\quad\Rightarrow\quad \sqrt{-g}\,R\;\to\;\sqrt{-g}\,\big(R-\Box_g\ln\alpha(x)\big),\qquad \alpha(x)>0 .
\label{eq:KK-weyl}
\end{equation}
The kinetic term in the action is
\begin{equation}
\frac{1}{2v^2}\sqrt{-g}\;g^{\mu\nu}\partial_\mu v\,\partial_\nu v .
\label{eq:KK-vkin}
\end{equation}
If $\alpha=v^b$, then
\begin{equation}
\begin{aligned}
\sqrt{-g}\,v\,R\;&\to\;\sqrt{-g}\;v\,\big(R-b\,\Box\ln v\big)\\
\end{aligned}
\label{eq:KK-weyl-v}
\end{equation}
So we rescale the metric by
\begin{equation}
g_{\mu\nu}\;\to\;\frac{r_h}{v^{1/2}}\;g_{\mu\nu}.
\label{eq:KK-rescale}
\end{equation}

which cancels the kinetic terms coming from term I and II. The action becomes ($v^2= \det \Phi$)
\begin{equation}
\begin{aligned}
S_2=&\frac{V_2}{16\pi G_4}\int_{\mathcal M} d^2x\,\sqrt{-g}\,\Big(v R-2 \Lambda\,\sqrt{v} r_h
-\frac{v}{4}\,g^{\mu\nu}\,\text{Tr}\!\big(M^{-1}\partial_\mu M\,M^{-1}\partial_\nu M\big)\Big)\\
&+\frac{V_2}{8\pi G_4}\int dx\,\sqrt{-h}\;vK\;+\;S_{\text{Maxwell}}
\end{aligned}
\label{eq:KK-S2}
\end{equation}

\begin{equation}
S_{\text{Maxwell}}=-\frac{V_2}{4G_4}\int d^2x\,\sqrt{-g}\,\sqrt{v}\,r_h\Big(F_{\mu\nu}F^{\mu\nu}\frac{v}{r_h^2}
+2\,(\partial_\mu A_a)(\partial_\nu A_b)\,\Phi^{ab}g^{\mu\nu}\frac{\sqrt{v}}{r_h} + \frac{Q_M}{v^2} )
\label{eq:Maxwell}
\end{equation}
\\
We restrict ourselves to the magnetic charge sector, therefore we can set $F_{\mu \nu}=0$ and $A_a=A_b=0$. \\

\begin{equation}\label{Eq: Maxwell}
S_{\text{Maxwell}}=-\frac{V_2}{4G_4}\int d^2x\,\sqrt{-g}\,\sqrt{v}\,r_h \frac{Q_M}{v^2}
\end{equation}

\\
%=====================================================================
\subsection{JT theory}
Now, to make contact with the JT theory, we expand $\Phi$. let
\begin{equation}
\Phi_{ij}=r_h^2\,(\mathbb{I}+\phi)^2_{ij},\qquad
v=\sqrt{\det\Phi}=r_h^2\,\det{}(\mathbb{I}+\phi).
\label{eq:KK-Phi-phi}
\end{equation}
 $\phi_{ij}(t,r)$ is a symmetric $2\times2$ matrix with components
$\phi_{xx},\phi_{xy}$  and $\phi_{yy}$.\\

First, we expand to linear order in $\phi$ to get
\begin{equation}
v=r_h^2\,(1+\Tr \phi)+\mathcal O(\phi^2).
\label{eq:KK-v-lin}
\end{equation}


Reduced action becomes

\begin{equation}
\begin{aligned}
S_2={}&\frac{V_2\,r_h^2}{16\pi G_4}\int_{\mathcal M}d^2x\,\sqrt{-g}\,\Big[(R-2 \Lambda)+\text{Tr}(\phi)\,\Big(R-\Lambda\Big)\Big]-\frac{V_2 r_h^2}{4G}\int d^2x\,\sqrt{-g}\,\, \frac{Q_M^2}{r_h^4}
\\&-\frac{V_2 r_h^2}{4G_4}\int d^2x\,\sqrt{-g}\,\, \frac{Q_M^2}{r_h^4}(\frac{\Tr(\phi)}{2}-  \phi_{xx}-\phi_{yy})+\frac{V_2\,r_h^2}{8\pi G_4}\int dx\,\sqrt{-h}\;K\,(1+\text{Tr}(\phi))\;   +\;\mathcal O(\phi^2)
\end{aligned}
\label{eq:KK-S2-lin}
\end{equation}
% NOTE: the value of \# was transcribed from a small handwritten note; please verify.

%=====================================================================
\subsection{Cubic vertices for tree level 4-point function}

Now we study the $\sim\phi^2$ and $\sim\phi^3$ terms for $\phi_{xy}$ tree level 4-point function. Since $\det\Phi=r_h^4\,[\det(\mathrm{I}+\phi)]^2$, We have $\sqrt{\det\Phi}=v=r_h^2\det(1+\phi)$ :
\begin{align}
\mathcal O(\phi^0):&\quad r_h^2 \label{eq:KK-sqrtdet-0}\\[4pt]
\mathcal O(\phi^1):&\quad r_h^2 \Tr (\phi) \label{eq:KK-sqrtdet-1}\\[4pt]
\mathcal O(\phi^2):&\quad r_h^2\left(\phi_{xx}\phi_{yy}-\phi_{xy}^2\right) \label{eq:KK-sqrtdet-2}\\[4pt]
\end{align}
Here, the $\mathcal O (\phi^2)$  term in the determinant gives mass to the shear mode via the gauge field term.\\

\textbf{Q=-$\frac{v}{4}g^{\mu\nu}\,\text{Tr}\!\big(M^{-1}\partial_\mu M\,M^{-1}\partial_\nu M\big)$ Term}

To simplify the analysis, we define, using \ref{Eq:term2}
\begin{equation}
    Q= \frac{-v}{4}g^{\mu\nu}\,\text{Tr}\!\big(\Phi^{-1}\partial_\mu \Phi\,\Phi^{-1}\partial_\nu \Phi\big) +\frac{g^{\mu \nu} }{2v}\partial_\mu v \partial_\nu v
\end{equation}
Let 
\begin{equation}
\begin{aligned}
Q_1&=\frac{-v}{4}g^{\mu\nu}\,\text{Tr}\!\big(\Phi^{-1}\partial_\mu \Phi\,\Phi^{-1}\partial_\nu \Phi\big) \\&=v g^{\mu\nu}\Big[\,\,\Tr(\partial_\mu\phi\,\partial_\nu\phi)\;-\;2\,\Tr(\phi\,\partial_\mu\phi\,\partial_\nu\phi)\Big]+\mathcal O(\phi^4),
\label{eq:KK-kin}
\end{aligned}
\end{equation}
and,
\begin{equation}
Q_2=\frac{g^{\mu \nu} }{2v}\partial_\mu v \partial_\nu v
\end{equation}

Now we study the terms in action(\ref{eq:KK-S2} and  \ref{Eq: Maxwell}) which are quadratic in $\phi_{xy}$.

\begin{equation}
\begin{aligned}
S_2\big|_ {\phi_{xy}^2}=\frac{V_2 r_h^2}{16\pi G_4}\int_\mathcal{M} d^2x \sqrt{-g}\big ((-R+\Lambda)\phi_{xy}^2-2g^{\mu \nu} \partial_\mu\phi_{xy}\partial_\nu\phi_{xy} \big)-\frac{V_2 r_h^2}{4G_4}\int_\mathcal{M} d^2x \sqrt{-g} \frac{Q_M^2}{r_h^5}
\end{aligned}
\label{eq:KK-kin-xy2}
\end{equation}
This is the kinetic term for the shear mode. This also gives us the cubic vertices with exchange of $h_{tt}, h_{tr} \text{ and } h_{rr}$. These exchanges exactly match with what we would have for a minimally coupled scalar. Note that $Q_2$ doesn't contribute any $\phi_{xy}^2$ term. Now, the point would be to show that other cubic vertex discussed below have contributions suppressed in $r_h$ compared to the $h_{tt}, h_{tr} \text{ and } h_{rr}$ exchanges.
\subsubsection*{(b) Cubic terms of the form $\phi_{xy}^2\times\phi$}
\begin{equation}
\begin{aligned}
K\Big|_{\phi_{xy}^2\phi}
=\;&g^{tt}\Big[-8\,(\phi_{xx}+\phi_{yy})(\partial_t\phi_{xy})^2
-16\,\phi_{xy}(\partial_t\phi_{xy})(\partial_t\phi_{xx}+\partial_t\phi_{yy})\Big]\\[4pt]
+\;&g^{tr}\Big[-16\,(\phi_{xx}+\phi_{yy})(\partial_t\phi_{xy})(\partial_r\phi_{xy})
-16\,\phi_{xy}(\partial_t\phi_{xy})(\partial_r\phi_{xx}+\partial_r\phi_{yy})\\
&\qquad-16\,\phi_{xy}(\partial_r\phi_{xy})(\partial_t\phi_{xx}+\partial_t\phi_{yy})\Big]\\[4pt]
+\;&g^{rr}\Big[-8\,(\phi_{xx}+\phi_{yy})(\partial_r\phi_{xy})^2
-16\,\phi_{xy}(\partial_r\phi_{xy})(\partial_r\phi_{xx}+\partial_r\phi_{yy})\Big]
\end{aligned}
\label{eq:KK-kin-xy2phi}
\end{equation}
Covariantly,
\begin{equation}
K\Big|_{\phi_{xy}^2\phi}=-8\,g^{\mu\nu}\Big[(\phi_{xx}+\phi_{yy})\,\partial_\mu\phi_{xy}\,\partial_\nu\phi_{xy}
+\partial_\mu\big(\phi_{xy}^2\big)\,\partial_\nu(\phi_{xx}+\phi_{yy})\Big].
\label{eq:KK-kin-xy2phi-cov}
\end{equation}
Combining all the terms, the relevant action for the $\phi_{xy}$ mode becomes:
\begin{equation}\label{phi action}
    S_{\phi_{xy}}=\frac{V_3 r_h^3}{16\pi G}\int_{\mathcal M} d^2x\,\sqrt{-g}\, \,\,\big(8 g^{\mu\nu}\partial_\mu\phi_{xy}\,\partial_\nu\phi_{xy}+ 4\pi F^2 \phi_{xy}^2 -8\,g^{\mu\nu}\Big[(\phi_{xx}+\phi_{yy})\,\partial_\mu\phi_{xy}\,\partial_\nu\phi_{xy}
+\partial_\mu\big(\phi_{xy}^2\big)\,\partial_\nu(\phi_{xx}+\phi_{yy})\Big] \big)
\end{equation}
Here, the cubic vertices in \ref{phi action} involve $\phi_{xx}$ and $\phi_{yy}$, which at leading order goes like $\mathcal{O}(\frac{r-r_h}{r_h})$,  whereas the $h_{tt}, h_{tr} \text{ and } h_{rr}$ exchanges are enhanced by $r_h$, as seen in \ref{epselon solution} or \\\ref{bdy side action}.  The dominant contribution to the tree-level four point function of $\phi_{xy}$ comes from \ref{eq:KK-kin-xy2}, which is $r_h^4$ enhanced, while the other vertices are suppressed by  $\frac{1}{r_h}$. This conclusively shows that shear mode behaves like a minimally coupled massive scalar in the near extremal blackhole background in the leading order in $r_h$. The mass of the scalar field is determined by $(\frac{Q}{r_h^3})^2$.

(Express the Q in terms of $r_h$ and write the mass in $L_2$ units)
%================================================================

\section{Bulk Action}

\label{sec:Bulk Action}



\appendix
\section{Details of Bulk calculation}
\label{sec: Appendix A}





\acknowledgments

This is the most common positions for acknowledgments. A macro is
available to maintain the same layout and spelling of the heading.

\paragraph{Note added.} This is also a good position for notes added
after the paper has been written.

% All figures and tables should be referenced in the text and should be
% placed on the page where they are first cited or in
% subsequent pages. Positioning them in the source file
% after the paragraph where you first reference them usually yield good
% results. See figure~\ref{fig:i} and table~\ref{tab:i} for layout examples. 
% Please note that a caption is mandatory  and it must be placed at the bottom of both figures and tables.

% \begin{figure}[htbp]
% \centering
% \includegraphics[width=.4\textwidth]{example-image-a}
% \qquad
% \includegraphics[width=.4\textwidth]{example-image-b}
% \caption{Always give a caption.\label{fig:i}}
% \end{figure}

% \begin{table}[htbp]
% \centering
% \begin{tabular}{lr|c}
% \hline
% x&y&x and y\\
% \hline
% a & b & a and b\\
% 1 & 2 & 1 and 2\\
% $\alpha$ & $\beta$ & $\alpha$ and $\beta$\\
% \hline
% \end{tabular}
% \caption{We prefer to have top and bottom borders around the tables.\label{tab:i}}
% \end{table}

% We discourage the use of inline figures (e.g. \texttt{wrapfigure}), as they may be
% difficult to position if the page layout changes.

% We suggest not to abbreviate: ``section'', ``appendix'', ``figure''
% and ``table'', but ``eq.'' and ``ref.'' are welcome. Also, please do
% not use \texttt{\textbackslash emph} or \texttt{\textbackslash it} for
% latin abbreviaitons: i.e., et al., e.g., vs., etc.


% \paragraph{Up to paragraphs.} We find that having more levels usually
% reduces the clarity of the article. Also, we strongly discourage the
% use of non-numbered sections (e.g.~\texttt{\textbackslash
%   subsubsection*}).  Please also consider the use of
% ``\texttt{\textbackslash texorpdfstring\{\}\{\}}'' to avoid warnings
% from the \texttt{hyperref} package when you have math in the section titles.

% Bibliography

%% [A] Recommended: using JHEP.bst file
%% \bibliographystyle{JHEP}
%% \bibliography{biblio.bib}

%% or
%% [B] Manual formatting (see below)
%% (i) We suggest to always provide author, title and journal data or doi:
%% in short all the informations that clearly identify a document.
%% (ii) please avoid comments such as "For a review'', "For some examples",
%% "and references therein" or move them in the text. In general, please leave only references in the bibliography and move all
%% accessory text in footnotes.
%% (iii) Also, please have only one work for each \bibitem.

% \begin{thebibliography}{99}

% \bibitem{a}
% Author,
% \emph{Title},
% \emph{J. Abbrev.} {\bf vol} (year) pg.

% \bibitem{b}
% Author,
% \emph{Title},
% arxiv:1234.5678.

% \bibitem{c}
% Author,
% \emph{Title},
% Publisher (year).

% \end{thebibliography}
\end{document}
